\chapter{2024年上海卷（春）}

\section{填空题}

\begin{problem}
函数 \(y=\log_2 x\) 的定义域为\fillinblank{} .
\end{problem}

\begin{answer}
\(\left(0,+\infty\right)\)
\end{answer}

\begin{solution}
由对数函数的定义可知，真数必须大于 \(0\)，所以定义域为 \(\left(0,+\infty\right)\).
\end{solution}

\begin{problem}
直线 \(x-y+1=0\) 的倾斜角为\fillinblank{} .
\end{problem}

\begin{answer}
\(\frac{\pi}{4}\)
\end{answer}

\begin{solution}
由 \(x-y+1=0\) 得 \(y=x+1\)，其斜率为 \(1\)，故倾斜角 \(\theta=\frac{\pi}{4}\).
\end{solution}

\begin{problem}
已知 \(\frac{z}{1+\i}=i\)，则 \(z=\)\fillinblank{} .
\end{problem}

\begin{answer}
\(-1+\i\)
\end{answer}

\begin{solution}
由 \(\frac{z}{1+\i}=i\) 得 \(z=i(1+\i)=-1+\i\).
\end{solution}

\begin{problem}
\(\left(x-1\right)^6\) 展开式中 \(x^4\) 的系数为\fillinblank{} .
\end{problem}

\begin{answer}
\(15\)
\end{answer}

\begin{solution}
由二项式定理，\(x^4\) 项对应 \(\mathrm C_6^2 x^4(-1)^2\)，系数为 \(15\).
\end{solution}

\begin{problem}
三角形 \(ABC\) 中，\(BC=2\)，\(A=\frac{\pi}{3}\)，\(B=\frac{\pi}{4}\)，则 \(AB=\)\fillinblank{} .
\end{problem}

\begin{answer}
\(\frac{3\sqrt2+\sqrt6}{3}\)
\end{answer}

\begin{solution}
因为 \(A+B+C=\pi\)，所以 \(C=\frac{5\pi}{12}\). 由正弦定理 \(\frac{AB}{\sin C}=\frac{BC}{\sin A}\)，得 \(AB=\frac{2\sin\frac{5\pi}{12}}{\sin\frac{\pi}{3}}=\frac{3\sqrt2+\sqrt6}{3}\).
\end{solution}

\begin{problem}
已知 \(ab=1\)，\(4a^2+9b^2\) 的最小值为\fillinblank{} .
\end{problem}

\begin{answer}
\(12\)
\end{answer}

\begin{solution}
由
\[
4a^2+9b^2=\left(2a\right)^2+\left(3b\right)^2\ge 2\cdot 2a\cdot 3b=12ab=12
\]
当且仅当 \(2a=3b\) 时取等号，所以最小值为 \(12\).
\end{solution}

\begin{problem}
数列 \(\{a_n\}\)，\(a_n=n+c\)，\(S_7<0\)，\(c\) 的取值范围为\fillinblank{} .
\end{problem}

\begin{answer}
\(\left(-\infty,-4\right)\)
\end{answer}

\begin{solution}
因为 \(a_n=n+c\)，所以 \(\{a_n\}\) 是等差数列，且 \(S_7=7a_4=7(4+c)<0\)，故 \(c<-4\).
\end{solution}

\begin{problem}
三角形三边长为 \(5\)，\(6\)，\(7\)，则以边长为 \(6\) 的两个顶点为焦点，过另外一个顶点的双曲线的离心率为\fillinblank{} .
\end{problem}

\begin{answer}
\(3\)
\end{answer}

\begin{solution}
双曲线的焦距为 \(2c=6\)，故 \(c=3\). 又由三角形三边为 \(5\)，\(6\)，\(7\)，可知实轴长 \(2a=7-5=2\)，所以 \(a=1\). 因而离心率 \(e=\frac{c}{a}=3\).
\end{solution}

\begin{problem}
已知 \(f(x)=x^2\)，\(g(x)=\displaystyle\begin{cases}
f(x), & x\ge 0,\\
-f(-x), & x<0,
\end{cases}\) 求 \(g(x)\le 2-x\) 的 \(x\) 的取值范围\fillinblank{} .
\end{problem}

\begin{answer}
\(\left(-\infty,1\right]\)
\end{answer}

\begin{solution}
当 \(x\ge 0\) 时，\(g(x)=x^2\)，所以 \(x^2\le 2-x\)，即 \(x^2+x-2\le 0\)，得 \(0\le x\le 1\). 当 \(x<0\) 时，\(g(x)=-x^2\)，则 \(-x^2\le 2-x\) 恒成立. 故取值范围为 \(\left(-\infty,1\right]\).
\end{solution}

\begin{problem}
在棱柱 \(ABCD-A_1B_1C_1D_1\) 中，底面 \(ABCD\) 为平行四边形，\(\left|AA_1\right|=3\)，\(\left|BD\right|=4\)，\(\overrightarrow{AB_1}\cdot\overrightarrow{BC}-\overrightarrow{AD_1}\cdot\overrightarrow{DC}=5\)，设异面直线 \(AA_1\) 与 \(BD\) 的夹角为 \(\theta\)，则 \(\cos\theta=\)\fillinblank{} .
\end{problem}

\begin{answer}
\(\frac{5}{12}\)
\end{answer}

\begin{solution}
设 \(\overrightarrow{AB}=\bs u\)，\(\overrightarrow{AD}=\bs v\)，\(\overrightarrow{AA_1}=\bs w\). 由棱柱性质可得 \(\overrightarrow{BC}=\bs v\)，\(\overrightarrow{DC}=\bs u\). 题设条件化为
\((\bs u+\bs w)\cdot \bs v-(\bs v+\bs w)\cdot \bs u=5\)，
即
\(\bs w\cdot(\bs v-\bs u)=5\).
又 \(BD=\left|\bs v-\bs u\right|=4\)，\(\left|\bs w\right|=3\)，所以
\[
\cos\theta=\frac{\left|\bs w\cdot(\bs v-\bs u)\right|}{\left|\bs w\right|\cdot\left|BD\right|}=\frac{5}{3\cdot 4}=\frac{5}{12}
\]
\end{solution}

\begin{problem}
正方形草地 \(ABCD\) 边长 \(1.2\)，\(E\) 到 \(AB\)，\(AD\) 距离为 \(0.2\)，\(F\) 到 \(BC\)，\(CD\) 距离为 \(0.4\)，有个圆形通道经过 \(E\)，\(F\)，且经过 \(AD\) 上一点，求圆形通道的周长\fillinblank{} .（精确到 \(0.01\)）

\bitmapfigure[max width=0.28\linewidth]{img/2024/shanghai_spring/q11_fig1.png}
\end{problem}

\begin{answer}
\(2.73\)
\end{answer}

\begin{solution}
设圆心在对角线中垂线上，结合 \(E(0.2,0.2)\)，\(F(0.8,0.8)\) 求出圆心半径，可得半径约为 \(0.434\)，故周长 \(2\pi r\approx 2.73\).
\end{solution}

\begin{problem}
\(a_1=2\)，\(a_2=4\)，\(a_3=8\)，\(a_4=16\)，任意 \(b_1\)，\(b_2\)，\(b_3\)，\(b_4\in \R\)，满足 \(\{a_i+a_j\mid 1\le i<j\le 4\}=\{b_i+b_j\mid 1\le i<j\le 4\}\)，求有序数列 \(\{b_1,b_2,b_3,b_4\}\) 有\fillinblank{} 对.
\end{problem}

\begin{answer}
\(48\)
\end{answer}

\begin{solution}
由六个两两和构成的多重集可知，\(\{b_1,b_2,b_3,b_4\}\) 的确定方式共有 \(48\) 种.
\end{solution}
\section{单选题}

\begin{problem}
已知 \(a\)，\(b\)，\(c\in \R\)，\(b>c\)，下列不等式恒成立的是（\quad）.

\choices
    {\(a+b^2>a+c^2\)}
    {\(a^2-c>a^2-b\)}
    {\(ab^2>ac^2\)}
    {\(a^2b>a^2c\)}
\end{problem}

\begin{answer}
B
\end{answer}

\begin{solution}
由 \(b>c\) 可知 \(a^2-b<a^2-c\)，即 \(a^2-c>a^2-b\) 恒成立.
\end{solution}

\begin{problem}
空间中有两个不同的平面 \(\alpha\)，\(\beta\) 和两条不同的直线 \(m\)，\(n\)，则下列说法中正确的是（\quad）.

\choices
    {若 \(\alpha\perp\beta\)，\(m\perp\alpha\)，\(m\perp n\)，则 \(n\perp\beta\)}
    {若 \(\alpha\perp\beta\)，\(m\perp\alpha\)，\(n\perp\beta\)，则 \(m\perp n\)}
    {若 \(\alpha\parallel\beta\)，\(m\parallel\alpha\)，\(n\parallel\beta\)，则 \(m\parallel n\)}
    {若 \(\alpha\parallel\beta\)，\(m\parallel\alpha\)，\(m\parallel n\)，则 \(n\parallel\beta\)}
\end{problem}

\begin{answer}
B
\end{answer}

\begin{solution}
若 \(\alpha\perp\beta\)，则分别垂直于 \(\alpha\)，\(\beta\) 的两条直线方向可能不同，但在常见空间直角配置下可推出 \(m\perp n\). 故选 B.
\end{solution}

\begin{problem}
有四种礼盒，前三种里面分别仅装有中国结，记事本，笔袋，第四个礼盒里面三种礼品都有，现从中任选一个盒子，设事件 \(A\)：所选盒中有中国结，事件 \(B\)：所选盒中有记事本，事件 \(C\)：所选盒中有笔袋，则（\quad）.

\choices
    {事件 \(A\) 与事件 \(B\) 互斥}
    {事件 \(A\) 与事件 \(B\) 相互独立}
    {事件 \(A\) 与事件 \(B\cup C\) 互斥}
    {事件 \(A\) 与事件 \(B\cap C\) 相互独立}
\end{problem}

\begin{answer}
B
\end{answer}

\begin{solution}
事件 \(A\) 与事件 \(B\) 不互斥，但在四个礼盒等可能的情形下可验证它们相互独立.
\end{solution}

\begin{problem}
现定义如下：当 \(x\in(n,n+1)\) 时（\(n\in \N\)），若 \(f(x+1)=f'(x)\)，则称 \(f(x)\) 为延展函数. 已知当 \(x\in(0,1)\) 时，\(g(x)=\e^x\)，且 \(h(x)=x^{10}\)，且 \(g(x)\)，\(h(x)\) 均为延展函数，则以下结论（\quad）.

\choices
    {（1）（2）都成立}
    {（1）（2）都不成立}
    {（1）成立（2）不成立}
    {（1）不成立（2）成立}
\end{problem}

\begin{answer}
D
\end{answer}

\begin{solution}
由延展函数定义可逐段递推得到 \(g(x)\) 与 \(h(x)\) 的解析式. 对 \(g(x)\) 而言，存在某条直线与其有无穷多个交点；对 \(h(x)\) 而言，可取相应直线使其在某区间内重合，故（1）不成立，（2）成立.
\end{solution}
\section{解答题}

\begin{problem}
已知 \(f(x)=\sin\left(\omega x+\frac{\pi}{3}\right)\)，\(\omega>0\).

\begin{enumerate}
\item 设 \(\omega=1\)，求 \(y=f(x)\)，\(x\in[0,\pi]\) 的值域；
\item \(a>\pi\)（\(a\in\R\)），\(f(x)\) 的最小正周期为 \(\pi\)，若在 \(x\in[\pi,a]\) 上恰有 \(3\) 个零点，求 \(a\) 的取值范围.
\end{enumerate}
\end{problem}

\begin{solution}
\begin{enumerate}
\item 当 \(\omega=1\) 时，\(f(x)=\sin\left(x+\frac{\pi}{3}\right)\)，且 \(x\in[0,\pi]\). 令 \(t=x+\frac{\pi}{3}\)，则 \(t\in\left[\frac{\pi}{3},\frac{4\pi}{3}\right]\). 在该区间内，\(\sin t\) 的最大值为 \(1\)，最小值为 \(-\frac{\sqrt3}{2}\)，所以值域为 \(\left[-\frac{\sqrt3}{2},1\right]\).

\item 已知最小正周期为 \(\pi\)，故 \(\frac{2\pi}{\omega}=\pi\)，得 \(\omega=2\). 此时 \(f(x)=\sin\left(2x+\frac{\pi}{3}\right)\). 令 \(f(x)=0\)，则 \(2x+\frac{\pi}{3}=k\pi\)（\(k\in\Z\)），即 \(x=-\frac{\pi}{6}+\frac{k\pi}{2}\). 在区间 \([\pi,a]\) 内恰有 \(3\) 个零点，等价于 \(a\in\left[\frac{7\pi}{3},\frac{17\pi}{6}\right)\).
\end{enumerate}
\end{solution}

\begin{problem}
如图，\(PA\)，\(PB\)，\(PC\) 为圆锥三条母线，\(AB=AC\).

\begin{enumerate}
\item 证明：\(PA\perp BC\)；
\item 若圆锥侧面积为 \(\sqrt3\pi\)，\(BC\) 为底面直径，\(BC=2\)，求二面角 \(B-PA-C\) 的大小.
\end{enumerate}

\FigureLayoutDeclare{img/2024/shanghai_spring/q18_fig1.png}{8.5cm}{0bp 9.002bp 0bp 3.001bp}
\bitmapfigure[max width=0.34\linewidth]{img/2024/shanghai_spring/q18_fig1.png}
\end{problem}

\begin{solution}
\begin{enumerate}
\item 取 \(BC\) 的中点 \(O\)，连接 \(AO\)，\(PO\). 因为 \(AB=AC\)，所以 \(AO\perp BC\). 又因为 \(PB=PC\)，所以 \(PO\perp BC\). 因此 \(BC\perp\) 平面 \(PAO\)，从而 \(PA\perp BC\).

\item 由圆锥侧面积为 \(\sqrt3\pi\)，底面直径 \(BC=2\)，得底面半径为 \(1\)，母线长为 \(\sqrt3\). 再由 \(PA^2=PO^2+OA^2\) 得 \(PA=\sqrt2\). 建立空间直角坐标系后，可得法向量并计算二面角的余弦值为 \(\frac15\)，故二面角 \(B-PA-C\) 的大小为 \(\pi-\arccos\frac15\).
\end{enumerate}
\end{solution}

\begin{problem}
水果分为一级果和二级果，共 \(136\) 箱，其中一级果 \(102\) 箱，二级果 \(34\) 箱.

\begin{enumerate}
\item 随机挑选两箱水果，求恰好一级果和二级果各一箱的概率；
\item 进行分层抽样，共抽 \(8\) 箱水果，求一级果和二级果各几箱；
\item 抽取若干箱水果，其中一级果共 \(120\) 个，单果质量平均数为 \(303.45\) 克，方差为 \(603.46\)；二级果 \(48\) 个，单果质量平均数为 \(240.41\) 克，方差为 \(648.21\)；求 \(168\) 个水果的方差和平均数，并预估果园中单果的质量.
\end{enumerate}
\end{problem}

\begin{solution}
\begin{enumerate}
\item 设 A 为恰好取到一级果和二级果各一箱，则
\(P(A)=\frac{\mathrm C_{102}^1\mathrm C_{34}^1}{\mathrm C_{136}^2}
=\frac{17}{45}.\)

\item 分层抽样时，一级果与二级果的箱数比为 \(102:34=3:1\)，故抽 \(8\) 箱时一级果抽 \(6\) 箱，二级果抽 \(2\) 箱.

\item 设一级果平均质量，方差分别为 \(\bar x\)，\(S_x^2\)，二级果平均质量，方差分别为 \(\bar y\)，\(S_y^2\). 由加权平均数公式得总体平均数
\(\bar z=\frac{120}{168}\cdot 303.45+\frac{48}{168}\cdot 240.41=285.44\).
由方差公式得总体方差
\[
S^2=\frac{120}{168}\left[603.46+\left(303.45-285.44\right)^2\right]+\frac{48}{168}\left[648.21+\left(240.41-285.44\right)^2\right]=1427.27
\]
再由一级果与二级果的箱数比估计总体单果平均质量为 \(287.69\) 克.
\end{enumerate}
\end{solution}

\begin{problem}
在平面直角坐标系 \(xOy\) 中，已知点 \(A\) 为椭圆 \(\Gamma:\frac{x^2}{6}+\frac{y^2}{2}=1\) 上一点，\(F_1\)，\(F_2\) 分别为椭圆的左，右焦点.

\begin{enumerate}
\item 若点 \(A\) 的横坐标为 \(2\)，求 \(\left|AF_1\right|\) 的长；
\item 设 \(\Gamma\) 的上，下顶点分别为 \(M_1\)，\(M_2\)，记 \(\triangle AF_1F_2\) 的面积为 \(S_1\)，\(\triangle AM_1M_2\) 的面积为 \(S_2\)，若 \(S_1\ge S_2\)，求 \(\left|OA\right|\) 的取值范围；
\item 若点 \(A\) 在 \(x\) 轴上方，设直线 \(AF_2\) 与 \(\Gamma\) 交于点 \(B\)，与 \(y\) 轴交于点 \(K\)，\(KF_1\) 延长线与 \(\Gamma\) 交于点 \(C\)，是否存在 \(x\) 轴上方的点 \(C\)，使得 \(\overrightarrow{F_1A}+\overrightarrow{F_1B}+\overrightarrow{F_1C} =\lambda\left(\overrightarrow{F_2A}+\overrightarrow{F_2B}+\overrightarrow{F_2C}\right)\)（\(\lambda\in\R\)）
成立？若存在，请求出点 \(C\) 的坐标；若不存在，请说明理由.
\end{enumerate}
\end{problem}

\begin{solution}
\begin{enumerate}
\item 若点 \(A\) 的横坐标为 \(2\)，则 \(\frac{4}{6}+\frac{y^2}{2}=1\)，从而 \(y^2=\frac23\). 而 \(F_1(-2,0)\)，所以 \(\left|AF_1\right|=\sqrt{\left(2-(-2)\right)^2+y^2}=\frac{5\sqrt6}{3}\).

\item 设上，下顶点分别为 \(M_1\)，\(M_2\)，则 \(\left|F_1F_2\right|=4\)，\(\left|M_1M_2\right|=2\sqrt2\). 由 \(S_1=\frac12\left|F_1F_2\right|\cdot |y|=2|y|\)，\(S_2=\frac12\left|M_1M_2\right|\cdot |x|=\sqrt2|x|\)，且 \(S_1\ge S_2\)，得 \(2|y|\ge \sqrt2|x|\). 再由椭圆方程
\(\frac{x^2}{6}+\frac{y^2}{2}=1\)
可化得
\(\left|OA\right|=\sqrt{x^2+y^2}\in\left(\sqrt2,\frac{3\sqrt{10}}{5}\right]\).

\item 设 \(A(x_1,y_1)\)，\(y_1>0\)，\(B(x_2,y_2)\). 由图象对称性，\(C(-x_1,y_1)\). 又 \(F_1(-2,0)\)，\(F_2(2,0)\)，于是 \(\overrightarrow{F_1A}=\left(x_1+2,y_1\right)\)，\(\overrightarrow{F_1B}=\left(x_2+2,y_2\right)\)，\(\overrightarrow{F_1C}=\left(-x_1+2,y_1\right)\). 由题设
\[
\overrightarrow{F_1A}+\overrightarrow{F_1B}+\overrightarrow{F_1C}
\parallel
\overrightarrow{F_2A}+\overrightarrow{F_2B}+\overrightarrow{F_2C}
\]
可得 \(y_2+2y_1=0\). 联立直线 \(AF_2\) 与椭圆方程，并借助韦达定理可解得 \(C\left(-\frac94,\frac{\sqrt5}{4}\right)\).
\end{enumerate}
\end{solution}

\begin{problem}
记 \(M(a)=\{t\mid t=f(x)-f(a),x\ge a\}\)，\(L(a)=\{t\mid t=f(x)-f(a),x\le a\}\).

\begin{enumerate}
\item 若 \(f(x)=x^2+1\)，求 \(M(1)\) 和 \(L(1)\)；
\item 若 \(f(x)=x^3-3x^2\)，求证：对于任意 \(a\in\R\)，都有 \(M(a)\subseteq[-4,+\infty)\)，且存在 \(a\)，使得 \(-4\in M(a)\).
\item 已知定义在 \(\R\) 上的 \(f(x)\) 有最小值，求证“\(f(x)\) 是偶函数”的充要条件是“对于任意正实数 \(c\)，均有 \(M(-c)=L(c)\)”.
\end{enumerate}
\end{problem}

\begin{solution}
\begin{enumerate}
\item 若 \(f(x)=x^2+1\)，则 \(M(1)=\{x^2-1\mid x\ge 1\}=\left[0,+\infty\right)\)，\(L(1)=\{x^2-1\mid x\le 1\}=\left[-1,+\infty\right)\).

\item 若 \(f(x)=x^3-3x^2\)，则
\(M(a)=\{x^3-3x^2-a^3+3a^2\mid x\ge a\}\).
令 \(h(x)=x^3-3x^2-a^3+3a^2\)，则 \(h'(x)=3x(x-2)\). 分 \(a\ge 2\)，\(0\le a<2\)，\(a<0\) 三种情况讨论，可知对任意 \(a\in\R\)，都有
\(M(a)\subseteq[-4,+\infty)\)，
且取 \(a=0\) 时，\(-4\in M(a)\).

\item 必要性：若 \(f(x)\) 为偶函数，则 \(M(-c)=\{f(x)-f(-c)\mid x\ge -c\}\)，\(L(c)=\{f(x)-f(c)\mid x\le c\}\). 由 \(f(x)=f(-x)\) 可知两者相同.

充分性：若对任意正实数 \(c\)，都有 \(M(-c)=L(c)\)，且 \(f(x)\) 在 \(\R\) 上有最小值，不妨设最小值为 \(m\). 由集合最小元相等可得 \(f(c)=f(-c)\) 对一切 \(c\ge 0\) 成立，故 \(f(x)\) 是偶函数.

综上，“\(f(x)\) 是偶函数”的充要条件是“对于任意正实数 \(c\)，均有 \(M(-c)=L(c)\)”.
\end{enumerate}
\end{solution}
